\chapter{Trigonometry in the Right Triangle}\label{ch:g9:trig}

In a right triangle, once you know one acute angle, the \emph{shape} of
the triangle is fixed --- only its size can vary. The ratios of its sides
therefore depend only on the angle: these ratios are the cosine, sine and
tangent. Together with the Pythagorean theorem, they let you compute
every side and every angle of a right triangle from very little
information.

\section{The Pythagorean theorem, again}

\begin{theorem}[Pythagoras]\label{thm:g9:trig:pythagoras}
In a triangle $ABC$ right-angled at $A$, the square of the hypotenuse is
the sum of the squares of the two legs:
\[
BC^2 = AB^2 + AC^2 .
\]
Conversely, if $BC^2 = AB^2 + AC^2$ in a triangle, then the triangle is
right-angled at $A$.
\end{theorem}
\admitted

\begin{example}\label{ex:g9:trig:pythagoras}
A right triangle has legs $AB = 5$ and $AC = 12$. Then
\[
BC^2 = 5^2 + 12^2 = 25 + 144 = 169,
\qquad\text{so}\quad BC = \sqrt{169} = 13 .
\]
To find a \emph{leg}: if $BC = 10$ and $AB = 6$, then
$AC^2 = BC^2 - AB^2 = 100 - 36 = 64$, so $AC = 8$.
\end{example}

\section{Cosine, sine, tangent}

In a right triangle, relative to an acute angle $\theta$, the three sides
have names: the \emph{hypotenuse} (opposite the right angle, the longest
side), the side \emph{opposite} to $\theta$, and the side
\emph{adjacent} to $\theta$ (the leg touching $\theta$).

\begin{definition}[Trigonometric ratios]\label{def:g9:trig:ratios}
For an acute angle $\theta$ of a right triangle:
\[
\cos\theta = \frac{\text{adjacent}}{\text{hypotenuse}},
\qquad
\sin\theta = \frac{\text{opposite}}{\text{hypotenuse}},
\qquad
\tan\theta = \frac{\text{opposite}}{\text{adjacent}} .
\]
These ratios\index{cosine}\index{sine}\index{tangent} depend only on
$\theta$, not on the size of the triangle (all right triangles with the
same acute angle are scalings of one another, by Thales).
\end{definition}

\begin{omfigure}
\begin{tikzpicture}[scale=1.15]
  \coordinate (A) at (0,0);
  \coordinate (B) at (4.2,0);
  \coordinate (C) at (4.2,2.2);
  \draw[thick] (A) -- (B) -- (C) -- cycle;
  \draw[thin] (B) ++(-0.28,0) -- ++(0,0.28) -- ++(0.28,0);
  \draw[omThm, ->] (0.9,0) arc (0:27.7:0.9);
  \node[omThm, font=\small] at (1.25,0.21) {$\theta$};
  \node[below, font=\small, omDef] at (2.1,0) {adjacent};
  \node[right, font=\small, omMeth] at (4.2,1.1) {opposite};
  \node[above left=-2pt, font=\small, omProp] at (2.2,1.25) {hypotenuse};
  \node[below left, font=\small] at (A) {$B$};
  \node[below right, font=\small] at (B) {$A$};
  \node[above right, font=\small] at (C) {$C$};
\end{tikzpicture}

{\small The three sides seen from the angle $\theta = \widehat B$: the
hypotenuse $[BC]$, the adjacent side $[BA]$, the opposite side $[AC]$.}
\end{omfigure}

\begin{proposition}[First properties]\label{prop:g9:trig:properties}
For every acute angle $\theta$:
\begin{enumerate}
  \item $0 < \cos\theta < 1$ and $0 < \sin\theta < 1$ (a leg is shorter
        than the hypotenuse);
  \item $\cos^2\theta + \sin^2\theta = 1$;
  \item $\tan\theta = \dfrac{\sin\theta}{\cos\theta}$.
\end{enumerate}
\end{proposition}

\begin{proof}
Write $a$ for the adjacent side, $o$ for the opposite side, $h$ for the
hypotenuse of a right triangle with angle $\theta$.

1. $0 < a < h$ and $0 < o < h$, so the two ratios are strictly between
$0$ and $1$.

2. By Pythagoras, $a^2 + o^2 = h^2$. Divide everything by $h^2$:
\[
\cos^2\theta + \sin^2\theta
= \frac{a^2}{h^2} + \frac{o^2}{h^2}
= \frac{a^2 + o^2}{h^2} = 1 .
\]

3. $\dfrac{\sin\theta}{\cos\theta} = \dfrac{o/h}{a/h} = \dfrac oa
= \tan\theta$.
\end{proof}

\begin{method}[Finding a side]\label{met:g9:trig:side}
To compute an unknown side when one side and one acute angle are known:
\begin{enumerate}
  \item mark the known angle and label the three sides (hypotenuse,
        opposite, adjacent) \emph{from that angle};
  \item choose the ratio (cos, sin or tan) that involves the known side
        and the wanted side;
  \item write the equation and solve it;
  \item sanity-check: the hypotenuse must come out longest.
\end{enumerate}
\end{method}

\begin{example}\label{ex:g9:trig:side}
In a right triangle, the hypotenuse measures $8$ and one angle is
$35^\circ$; compute the side opposite to it. The ratio involving the
opposite side and the hypotenuse is the sine:
\[
\sin 35^\circ = \frac{\text{opposite}}{8}
\quad\Longrightarrow\quad
\text{opposite} = 8 \sin 35^\circ \approx 8 \times 0.574 \approx 4.6 .
\]
For the adjacent side, use the cosine:
$8\cos 35^\circ \approx 8 \times 0.819 \approx 6.6$. Check with
Pythagoras: $4.6^2 + 6.6^2 \approx 21 + 43 \approx 64 = 8^2$.
\end{example}

\begin{method}[Finding an angle]\label{met:g9:trig:angle}
When two sides are known, compute the ratio they form, identify it as a
cosine, sine or tangent of the unknown angle, and apply the calculator's
inverse function ($\cos^{-1}$, $\sin^{-1}$ or $\tan^{-1}$) to the ratio.
\end{method}

\begin{example}\label{ex:g9:trig:angle}
A ladder $5$ m long leans against a wall, its foot $1.4$ m from the
wall. The angle $\theta$ between the ladder and the ground satisfies
\[
\cos\theta = \frac{1.4}{5} = 0.28,
\qquad\text{so}\quad
\theta = \cos^{-1}(0.28) \approx 73.7^\circ .
\]
\end{example}

\begin{omfigure}
\begin{tikzpicture}[scale=0.85]
  \draw[thick] (0,0) -- (1.4,0);
  \draw[thick] (1.4,0) -- (1.4,4.0);
  \draw[omDef, very thick] (0,0) -- (1.4,4.0);
  \draw[thin] (1.4,0) ++(-0.25,0) -- ++(0,0.25) -- ++(0.25,0);
  \draw[omThm, ->] (0.75,0) arc (0:70.7:0.75);
  \node[omThm, font=\small] at (1.15,0.55) {$\theta$};
  \node[below, font=\small] at (0.7,0) {$1.4$ m};
  \node[omDef, font=\small, left] at (0.62,2.1) {$5$ m};
  \draw[gray, thin] (-0.7,0) -- (2.3,0);
\end{tikzpicture}

{\small The ladder problem: the adjacent side ($1.4$ m) and the
hypotenuse ($5$ m) are known, so the cosine gives the angle.}
\end{omfigure}

\begin{example}[Special angles]\label{ex:g9:trig:special}
Two triangles give exact values. Half a square of side 1 (a right
isosceles triangle) has hypotenuse $\sqrt2$ and angles of $45^\circ$:
\[
\cos 45^\circ = \sin 45^\circ = \frac{1}{\sqrt2} = \frac{\sqrt2}{2},
\qquad \tan 45^\circ = 1 .
\]
Half an equilateral triangle of side 1 gives the angles $30^\circ$ and
$60^\circ$:
\[
\cos 60^\circ = \sin 30^\circ = \frac12,
\qquad
\cos 30^\circ = \sin 60^\circ = \frac{\sqrt3}{2}.
\]
\end{example}

\section{Exercises}

\begin{exercise}[$\star$]\label{exo:g9:trig:1}
A right triangle has legs $9$ and $12$. Compute its hypotenuse. Another
has hypotenuse $17$ and one leg $8$: compute the other leg.
\end{exercise}

\begin{exercise}[$\star$]\label{exo:g9:trig:2}
A triangle has sides $7$, $24$ and $25$. Is it right-angled? Same
question for sides $5$, $6$ and $8$.
\end{exercise}

\begin{exercise}[$\star$]\label{exo:g9:trig:3}
In a triangle $DEF$ right-angled at $D$, with the angle at $E$ noted
$\theta$: which side is the hypotenuse? Which side is opposite to
$\theta$? Write $\cos\theta$, $\sin\theta$ and $\tan\theta$ as ratios of
the sides $DE$, $DF$, $EF$.
\end{exercise}

\begin{exercise}[$\star$]\label{exo:g9:trig:4}
In a right triangle, the hypotenuse measures $10$ and one acute angle is
$28^\circ$. Compute the two legs (round to the tenth;
$\cos 28^\circ \approx 0.883$, $\sin 28^\circ \approx 0.469$).
\end{exercise}

\begin{exercise}[$\star$]\label{exo:g9:trig:5}
In a right triangle, the leg adjacent to the angle $\theta$ measures $6$
and the opposite leg measures $4.5$. Compute $\tan\theta$, then $\theta$
(round to the degree; $\tan^{-1}(0.75) \approx 36.9^\circ$).
\end{exercise}

\begin{exercise}[$\star\star$]\label{exo:g9:trig:6}
A drone flies at an altitude of $120$ m. From an observer on the ground,
it is seen at an angle of elevation of $32^\circ$. At what horizontal
distance from the observer is the drone
($\tan 32^\circ \approx 0.625$)?
\end{exercise}

\begin{exercise}[$\star\star$]\label{exo:g9:trig:7}
An access ramp must rise $1.2$ m with an angle of at most $5^\circ$ to
the horizontal. What minimum length along the slope must the ramp have
($\sin 5^\circ \approx 0.0872$)? Round up to the nearest tenth of a
meter.
\end{exercise}

\begin{exercise}[$\star\star$]\label{exo:g9:trig:8}
Using \cref{prop:g9:trig:properties}: an acute angle $\theta$ satisfies
$\cos\theta = 0.6$. Compute $\sin\theta$ without finding $\theta$, then
$\tan\theta$.
\end{exercise}

\begin{exercise}[$\star\star\star$]\label{exo:g9:trig:9}
Justify the exact values of \cref{ex:g9:trig:special}: in the right
isosceles triangle with legs $1$, compute the hypotenuse; in the
equilateral triangle of side $1$, compute the height, and deduce the
cosine and sine of $30^\circ$ and $60^\circ$.
\end{exercise}
